% !TEX root = ../monograph.tex
\begin{proof}[Proof of Proposition~\ref{ch18:replacement-input}]
\label{ch18:raw-replacement}
The lattice estimate in Proposition~\ref{ch18:lattice-replacement}
was proved before phase decomposition. We use its uniformity in
replica number to obtain the algebra-valued statement for the
reconstructions $H_D$ fixed above.

Let $F$ be an interior lattice difference with chessboard seminorm
at most $B$, and let its exterior factors have interior seminorms
at most $B_1,\ldots,B_q$. Put copies of the complete difference
and exterior factors in separated cells before taking limits.
The repeated-word identity and chessboard inequality give, for its
compressed limit $Z$ and each integer $n\ge1$,
\[
 \int_\Xi|\xi(Z)|^{2n}\,d\nu(\xi)
                   \le(B\prod_{l=1}^qB_l)^{2n}.
\]
The seminorm bounds $B,B_1,\ldots,B_q$ are independent of $n$.
For each fixed $n$, the cutoff and the separation of the copies may
need to be increased.
This uses the complete signed density, with every term involving at least two insertions,
from the proof of \eqref{ch18:lattice-estimate}.
Equation~\eqref{ch15:moment-norm} then gives
$\|Z\|\le B\prod_lB_l$. Positivity here belongs to the limiting
commutative algebra and the original Wilson law, not to an
individual signed effective activity.

Apply this observation to the auxiliary field $A_t$ in the proof
of Proposition~\ref{ch18:mixed-bound}. On grid configurations its
replacement error has norm with exponent
$Z_0-(t+1/16)(\theta-\theta_-)$, where
$Z_0=W\theta_++4n_{\rm cl}(\theta_+-\theta)$.
Fixed exterior factors cost a bounded factor. The graded quotient
bound gives exponent $Z_0-i\theta$ for a degree-$i$ component of
$A_t$; all the finitely many fields used in its reconstruction
have uniformly bounded compressed exterior words. Testing the
same error against the fixed probes gives the probe coordinate
errors. With $t$ chosen as in that proof, the omitted degrees
and the low-coordinate errors therefore give
\[
 \|M(wv)-M((H_D(\tau_{-o}w))(o)v)\|
                \preccurlyeq e^{[Z_0-(D+1)\theta]T}.
\]
There is no spectator-dependent choice of $H_D$: the only
geometry-dependent auxiliary field has disappeared from this
inequality.

The derivative descent of Proposition~\ref{ch18:mixed-bound}
works in this norm. Equivalently, apply it to every character,
using the same sharp chessboard seminorm bounds, and take their supremum.
The crude and grid constants are uniform in the character.
A positional derivative changes its input degree by one; a finite
descending induction removes the grid restriction, also for any
fixed list of differentiated fields. Letting the two resolution
heights tend to one yields
\[
 \|M(wv)-M((H_D(\tau_{-o}w))(o)v)\|
                      \le C r^{D+1-W-\varepsilon}.
\]
The coefficient assertion is exactly the corresponding fixed-probe
estimate in Proposition~\ref{ch18:mixed-bound}. Compact label
sets, conjugation, and the finite probe families do not change
its constants. For their rotation orbits the same uniform estimates
apply on the compact set of rotated labels and positions; choose
$r$ so all these probes remain separated from the cluster.

Finally apply the vacuum state to $(Z_D^*Z_D)^n$. Its norm is
at most $\|Z_D\|^{2n}$, giving
\eqref{ch18:replacement-moments} with the same $C$ and the same
small-$r$ threshold for every $n$. Thus the order of quantifiers
is obtained from the norm estimate, rather than from an
uncontrolled sequence of separate moment constants.
\end{proof}
